\scanpage{106}
Donc on peut supposer $E=\complement V$, $V$ \deleted{est} ouvert rétrocompact,
\deleted{et par suite on aura $\unsure{E-F}=F=\complement\unsure{U}$ \unsure{le cas}}, Il suffit de prouver \inserted{\unsure{le puis}} $E-F$ const. dans $X$, i.e.
\[
X-(E-F)=V\cup F\text{ rétrocompact dans }X,
\]
ce qui est le cas car $V$ et $F$ le sont évidemment.

\margin{\diagram{\begin{tikzcd}[ampersand replacement=\&,row sep=small]
\complement V\arrow[d,no head]\vphantom{F}\\[-1ex]
\unsure{F}\arrow[d,no head]\\[-1ex]
\unsure{\cdot}
\end{tikzcd}}}

Le cas ``local'' résulte aussitôt du cas global.

\uline{Corollaire} 1\quad Supposons que de plus, \unsure{$\forall$} partie $U_i$ soit \uline{quasi compacte} \deleted{\unsure{fermée}} \uline{pour la topologie constructible} [donc aussi \deleted{pour la top} compacts pour la top constructible, dans le cas où \deleted{\unsure{la top}} $X$ satisfait $(T_0)$], \inserted{i.e. \unsure{satisfaisant} à la condition: $\unsure{\#}\ \unsure{---}\ \unsure{!}$} Alors

iv) Les $U_i^{\mathrm{cons}}$ sont quasi-compacts, \deleted{et $X$ \unsure{est donc}}
\deleted{\unsure{ouvert} \unsure{quasi} \unsure{cpt} une base formée d'ouverts}
\deleted{\unsure{pour} $X$ \unsure{admet} une base formée d'ouverts}

(v) Soit $E$ une partie \deleted{loc. constructible de}\inserted{fermée de} $X^{\mathrm{cons}}$ [p.ex. une partie loc. constructible de $X$] alors
\[
\begin{aligned}
E\ \deleted{\text{compact}}\ \text{quasi-compact}\ \inserted{\unsure{\text{(fort. q.cpt)}}}\ \text{dans }X^{\mathrm{cons}}
&\Longleftrightarrow E\text{ quasi-compact dans }X\\
&\Longleftrightarrow E\subset\text{réunion finie des }U_i.
\end{aligned}
\]

(vi) Pour que $X^{\mathrm{cons}}$ soit loc. cpt [cpt] il f. et s. que $X$ satisfasse $(T_0)$ [resp. $(T_0)$, et soit q. cpt].

\margin{N.B. Les hypothèses faites impliquent que \unsure{tt} fermé irred. de $X$ \unsure{admet} un \deleted{\unsure{pt}} pt générique (\unsure{compte tenu de} $T_0$ \unsure{est réciproque|et réciproque}\unsure{---})}

\margin{\figure{\begin{tikzpicture}[x=1cm,y=1cm]
\draw (0,0)--(0.12,-1.7);
\draw (0.2,-2.0)--(0.23,-2.7);
\draw (0.3,-3.1)--(0.38,-5.2)--(0.58,-5.22);
\end{tikzpicture}}}

(iv) Résulte de (i) trivialement. Évident
\[
E^{\mathrm{cons}}\text{ q. cpt}\Longrightarrow E\text{ q. cpt en vertu de (i)}
\]
$\Longrightarrow E\subset$ union finie \inserted{$V$} des $U_i$ par définition de q. cpt, donc $E$ fermé dans $V^{\mathrm{cons}}$ [pour la top. induite par $X^{\mathrm{cons}}$] qui est aussi $\bigcup U_i^{\mathrm{cons}}$, et les $U_i^{\mathrm{cons}}$ sont q. cpt, donc $E$ l'est, C.Q.

\margin{(vi) En résulte trivialement}
\margin{105}
